\documentclass[a4paper,10pt]{article}
\usepackage[utf8]{inputenc}
\usepackage[T1]{fontenc}
\usepackage[landscape, margin=1.5cm]{geometry}
\usepackage{array}
\usepackage{tabularx}
\usepackage{hyperref}
% xcolor con la opcion table: es la que define \rowcolors y carga colortbl.
% Sin ella el documento no compila (el archivo original tenia este fallo).
\usepackage[table]{xcolor}

\begin{document}

\section*{Comparativa de Modelos de IA Generativa (Versiones Gratuitas)}

Material complementario del \textbf{capítulo 4} de \textit{Inteligencia artificial generativa en la Educación}.

\textbf{Verificado el 10 de septiembre de 2026} contra la documentación oficial de cada proveedor.
Las filas marcadas con $\dagger$ se comprobaron en esa fecha; las demás proceden de revisiones
anteriores y \textbf{no pudieron reverificarse} porque sus proveedores no publican esta información
de forma accesible. Las celdas con «---» indican dato no confirmado.

\textbf{Este documento caduca.} Los planes gratuitos y sus límites cambian en cuestión de meses.
Antes de recomendar una herramienta a un estudiante, contrasta con la página oficial del proveedor
siguiendo la lista de verificación del repositorio.

\vspace{0.4cm}

\rowcolors{2}{gray!10}{white}
\begin{tabularx}{\textwidth}{|>{\raggedright\arraybackslash}X|
>{\centering\arraybackslash}m{2.3cm}|
>{\centering\arraybackslash}m{2.3cm}|
>{\centering\arraybackslash}m{2.3cm}|
>{\centering\arraybackslash}m{2.3cm}|
>{\centering\arraybackslash}m{2.1cm}|
>{\centering\arraybackslash}m{2.1cm}|
>{\raggedright\arraybackslash}X|}
\hline
\rowcolor{blue!40}
\textbf{Modelo} & \textbf{Entrada imágenes} & \textbf{Genera imágenes} & \textbf{Deep search} & \textbf{Ejecuta código} & \textbf{Búsqueda web} & \textbf{Datos usados para entrenar} & \textbf{Enlace} \\
\hline
ChatGPT $\dagger$ & Sí, con límites & Sí, con límites & Sí (Deep Research) & Sí, con límites & \textbf{Sí} & Sí, por defecto & \href{https://chatgpt.com}{chatgpt.com} \\
\hline
Claude $\dagger$ & Sí & No & No, pero con búsqueda web & \textbf{Sí} & \textbf{Sí} & No por defecto & \href{https://claude.ai}{claude.ai} \\
\hline
Gemini $\dagger$ & Sí & Sí, genera y edita & Sí (Deep Research) & Sí & Sí & Sí, por defecto & \href{https://gemini.google.com}{gemini.google.com} \\
\hline
Mistral (Vibe) $\dagger$ & Sí & Sí & --- & Sí, sesiones limitadas & Sí, limitada & Sí, con exclusión & \href{https://mistral.ai}{mistral.ai} \\
\hline
Grok $\dagger$ & Sí & Limitada & Sí (DeepSearch) & --- & Sí & No por defecto & \href{https://grok.com}{grok.com} \\
\hline
DeepSeek & Sí & --- & Sí, chat con web & --- & --- & Sí, por defecto & \href{https://chat.deepseek.com}{chat.deepseek.com} \\
\hline
Qwen Studio & Sí, imagen y vídeo & Sí & --- & --- & No documentada & Documentación limitada & \href{https://chat.qwen.ai}{chat.qwen.ai} \\
\hline
Kimi & Sí & --- & Sí (Researcher) & --- & --- & Sin exclusión clara & \href{https://www.kimi.com}{kimi.com} \\
\hline
Meta AI & Sí, fotos en chat & Sí, genera y edita & No & No & Sí & Sí, exclusión compleja & \href{https://www.meta.ai}{meta.ai} \\
\hline
\end{tabularx}

\vspace{0.5cm}

\noindent\textbf{Tres advertencias.}

\begin{enumerate}
  \item \textbf{Gratuito ya no significa lo mismo en todas partes.} OpenAI anunció que empezará a
  probar publicidad en sus planes gratuito y Go; los planes Pro, Business y Enterprise no la
  incluyen. Y en casi todos los proveedores el plan gratuito es justamente aquel en el que los datos
  se usan para entrenamiento por defecto.
  \item \textbf{Existen planes intermedios de bajo costo.} ChatGPT Go y Google AI Plus se dirigen al
  estudiante que agota el plan gratuito sin poder pagar una suscripción completa.
  \item \textbf{Mistral renombró su asistente a Vibe.} «Le Chat» sobrevive solo como nombre de la
  aplicación móvil en las tiendas.
\end{enumerate}

\vspace{0.3cm}

\noindent\textbf{Sobre las cuatro filas sin $\dagger$.} DeepSeek, Qwen, Kimi y Meta AI no pudieron
verificarse: Meta AI impide la consulta automatizada de su sitio, y las otras tres no publican
precios ni capacidades del servicio de consumo de forma accesible. Preferimos dejar el hueco visible
antes que rellenarlo con datos de terceros sin comprobar. Esa opacidad es, en sí misma, parte del
riesgo de adoptar una herramienta en una institución.

\end{document}
